\documentclass[11pt,a4paper]{article}

\usepackage[utf8]{inputenc}
\usepackage[T1]{fontenc}
\usepackage[spanish,es-tabla]{babel}
\usepackage[margin=2.5cm]{geometry}
\usepackage{amsmath}
\usepackage{booktabs}
\usepackage{array}
\usepackage{xcolor}
\usepackage{listings}
\usepackage{hyperref}
\usepackage{enumitem}

\hypersetup{colorlinks=true, linkcolor=blue!50!black, urlcolor=blue!50!black}

% ---------- Estilo para código Java ----------
\lstdefinestyle{javastyle}{
  language=Java,
  basicstyle=\ttfamily\footnotesize,
  keywordstyle=\color{blue!60!black}\bfseries,
  commentstyle=\color{green!40!black}\itshape,
  stringstyle=\color{red!60!black},
  numbers=left,
  numberstyle=\tiny\color{gray},
  frame=single,
  rulecolor=\color{gray!60},
  breaklines=true,
  showstringspaces=false,
  tabsize=4,
  literate={á}{{\'a}}1 {é}{{\'e}}1 {í}{{\'i}}1 {ó}{{\'o}}1 {ú}{{\'u}}1
           {Á}{{\'A}}1 {É}{{\'E}}1 {Í}{{\'I}}1 {Ó}{{\'O}}1 {Ú}{{\'U}}1
           {ñ}{{\~n}}1 {¿}{{?`}}1 {¡}{{!`}}1
}
\lstset{style=javastyle}

\title{\textbf{Actividad evaluativa \#1}\\[4pt]
\large Centro de monitoreo de accesos\\[4pt]
\large Reporte técnico del programa \texttt{Main.java}}
\author{Adrià Jolis \\ 2.º curso -- Ciberseguridad -- Java\\[6pt]
\textbf{Estudiante:} \underline{\hspace{6cm}}}
\date{23 de octubre de 2026}

\begin{document}

\maketitle
\tableofcontents
\newpage

% =====================================================
\section{Introducción y objetivo}
Un equipo de ciberseguridad necesita un programa que analice la actividad de
\textbf{dos servidores independientes}. Cada servidor recibe intentos de
autenticación y el sistema debe:
\begin{enumerate}[nosep]
  \item validar los registros introducidos,
  \item clasificar el riesgo de cada servidor (\texttt{BLOQUEADO}, \texttt{ALERTA}, \texttt{NORMAL} o \texttt{DATOS INVALIDOS}),
  \item determinar cuál de los dos requiere atención prioritaria.
\end{enumerate}

El programa se ha implementado en Java en un único archivo \texttt{Main.java}
que contiene la clase \texttt{Servidor} y la clase pública \texttt{Main}. Se han
utilizado únicamente los conocimientos permitidos: variables, tipos básicos,
métodos de \texttt{String}, \texttt{Scanner}, operadores lógicos y aritméticos,
condicionales, métodos, clases, objetos y \texttt{printf()}.

% =====================================================
\section{Diseño orientado a objetos}

\subsection{Clase \texttt{Servidor}}
La clase modela un servidor mediante cinco atributos de instancia:

\begin{table}[h!]
\centering
\begin{tabular}{@{}lll@{}}
\toprule
\textbf{Atributo} & \textbf{Tipo} & \textbf{Significado} \\
\midrule
\texttt{nombre}           & \texttt{String} & Nombre del servidor \\
\texttt{origen}           & \texttt{String} & \texttt{interno} o \texttt{externo} \\
\texttt{critico}          & \texttt{String} & \texttt{si} o \texttt{no} \\
\texttt{intentosTotales}  & \texttt{int}    & Total de intentos de autenticación \\
\texttt{intentosFallidos} & \texttt{int}    & Intentos fallidos \\
\bottomrule
\end{tabular}
\caption{Atributos de instancia de la clase \texttt{Servidor}.}
\end{table}

Los métodos exigidos por el enunciado son:

\begin{table}[h!]
\centering
\begin{tabular}{@{}>{\ttfamily}l p{9.5cm}@{}}
\toprule
\textbf{Método} & \textbf{Responsabilidad} \\
\midrule
boolean datosValidos()  & Comprueba todas las reglas de validación. \\
String evaluarEstado()  & Valida y devuelve el estado del servidor. \\
double porcentajeFallos() & Calcula el porcentaje de fallos; solo se usa con datos válidos. \\
void mostrarResumen()   & Imprime nombre, fallos/total, porcentaje (o \texttt{N/D}) y estado. \\
\bottomrule
\end{tabular}
\caption{Métodos exigidos por el enunciado.}
\end{table}

Además, se añadió un método auxiliar, \texttt{int nivelPrioridad()}, que
convierte el estado en un número (\texttt{BLOQUEADO}=3, \texttt{ALERTA}=2,
\texttt{NORMAL}=1, \texttt{DATOS INVALIDOS}=0) para poder comparar dos
servidores con operadores relacionales sencillos.

\subsection{Clase \texttt{Main}}
Contiene el método \texttt{main()} y el método estático auxiliar
\texttt{leerServidor(Scanner, int)}, que lee por consola los cinco datos de un
servidor y devuelve un objeto \texttt{Servidor} nuevo. En \texttt{main()} se
llama dos veces a este método, por lo que se obtienen \textbf{dos objetos
distintos e independientes} (\texttt{a} y \texttt{b}), cada uno con su propia
zona de memoria y sus propios valores.

% =====================================================
\section{Entrada de datos y tratamiento del texto}
Todos los datos se leen con \texttt{sc.nextLine()}. Los campos numéricos se
leen como texto y se convierten con \texttt{Integer.parseInt(...)}, aplicando
\texttt{trim()} antes:

\begin{lstlisting}
System.out.print("Intentos totales: ");
s.intentosTotales = Integer.parseInt(sc.nextLine().trim());
\end{lstlisting}

Esta decisión evita el problema clásico del \textbf{salto de línea pendiente}
que deja \texttt{nextInt()} en el búfer y que provocaría que la siguiente
lectura con \texttt{nextLine()} devolviera una cadena vacía.

Para las comparaciones de texto se usa \texttt{trim()} (elimina espacios al
principio y al final) junto con \texttt{equalsIgnoreCase()} (ignora
mayúsculas), tal como pide el enunciado. Así, entradas como
\texttt{"\ EXTERNO\ "} o \texttt{"\ SI\ "} se interpretan correctamente.

% =====================================================
\section{Validación de datos}
El método \texttt{datosValidos()} devuelve \texttt{false} en cuanto una regla
falla y \texttt{true} solo si se cumplen todas:

\begin{enumerate}[nosep]
  \item \textbf{Nombre:} no es \texttt{null} y \texttt{nombre.trim().isEmpty()} es falso (al menos un carácter que no sea espacio).
  \item \textbf{Origen y criticidad:} \texttt{origen} es \texttt{interno} o \texttt{externo}; \texttt{critico} es \texttt{si} o \texttt{no}.
  \item \textbf{Intentos:} \texttt{intentosTotales > 0} y \texttt{0 <= intentosFallidos <= intentosTotales} (extremos incluidos).
\end{enumerate}

\begin{lstlisting}
if (intentosTotales <= 0) {
    return false;
}
if (intentosFallidos < 0 || intentosFallidos > intentosTotales) {
    return false;
}
\end{lstlisting}

% =====================================================
\section{Clasificación del riesgo}
\texttt{evaluarEstado()} primero llama a \texttt{datosValidos()}; si falla,
devuelve \texttt{DATOS INVALIDOS} sin clasificar. Si los datos son válidos, se
normalizan una sola vez dos booleanos (\texttt{esExterno}, \texttt{esCritico})
y se aplican las reglas \textbf{en orden de prioridad}:

\begin{table}[h!]
\centering
\begin{tabular}{@{}c l p{9cm}@{}}
\toprule
\textbf{Orden} & \textbf{Estado} & \textbf{Condición} \\
\midrule
1 & \texttt{BLOQUEADO} & 5 o más fallos; o crítico, externo y al menos 3 fallos. \\
2 & \texttt{ALERTA}    & 3 o 4 fallos; o externo con al menos 2 fallos (solo si no es \texttt{BLOQUEADO}). \\
3 & \texttt{NORMAL}    & Datos válidos que no cumplen las condiciones anteriores. \\
\bottomrule
\end{tabular}
\caption{Reglas de clasificación.}
\end{table}

El orden de los \texttt{if} garantiza la prioridad: como \texttt{BLOQUEADO} se
comprueba primero y cada condición termina con \texttt{return}, la regla de
\texttt{ALERTA} solo se evalúa cuando el servidor no está bloqueado.

\begin{lstlisting}
if (intentosFallidos >= 5
        || (esCritico && esExterno && intentosFallidos >= 3)) {
    return "BLOQUEADO";
}
if (intentosFallidos >= 3
        || (esExterno && intentosFallidos >= 2)) {
    return "ALERTA";
}
return "NORMAL";
\end{lstlisting}

% =====================================================
\section{Porcentaje de fallos y resumen}
El porcentaje se calcula como
\[
  \text{porcentaje} = \frac{\text{intentosFallidos} \times 100}{\text{intentosTotales}}
\]
usando un \textit{cast} a \texttt{double} para evitar la división entera, que
truncaría el resultado:

\begin{lstlisting}
return (double) intentosFallidos * 100 / intentosTotales;
\end{lstlisting}

\texttt{mostrarResumen()} imprime con \texttt{printf} el formato
\texttt{nombre | fallidos/total | porcentaje | estado}. Con datos válidos se
muestra el porcentaje con un decimal (se usa \texttt{Locale.US} para que el
separador decimal sea siempre un punto); con datos inválidos se escribe
\texttt{N/D} y \textbf{no se efectúa ninguna división}.

% =====================================================
\section{Comparación entre los dos servidores}
En \texttt{main()} la comparación se resuelve mediante condiciones y el método
\texttt{nivelPrioridad()}, sin respuestas fijas:

\begin{enumerate}[nosep]
  \item Si alguno de los dos tiene datos inválidos $\rightarrow$ \texttt{COMPARACION NO DISPONIBLE}.
  \item Si ambos tienen el mismo nivel $\rightarrow$ \texttt{IGUAL PRIORIDAD}.
  \item En otro caso $\rightarrow$ \texttt{ATENCION PRIORITARIA:} seguido del nombre del servidor con mayor nivel.
\end{enumerate}

\begin{lstlisting}
if (!a.datosValidos() || !b.datosValidos()) {
    System.out.println("COMPARACION NO DISPONIBLE");
} else if (a.nivelPrioridad() == b.nivelPrioridad()) {
    System.out.println("IGUAL PRIORIDAD");
} else if (a.nivelPrioridad() > b.nivelPrioridad()) {
    System.out.println("ATENCION PRIORITARIA: " + a.nombre.trim());
} else {
    System.out.println("ATENCION PRIORITARIA: " + b.nombre.trim());
}
\end{lstlisting}

% =====================================================
\section{Pruebas}

\subsection{Ejemplo de ejecución del enunciado}
Entradas: \texttt{VPN / EXTERNO / SI / 10 / 3} y \texttt{CORREO / interno / no / 12 / 3}.

\begin{lstlisting}[numbers=none]
========== REPORTE DE SEGURIDAD ==========
VPN | 3/10 | 30.0% | BLOQUEADO
CORREO | 3/12 | 25.0% | ALERTA
ATENCION PRIORITARIA: VPN
\end{lstlisting}

\noindent Razonamiento: \texttt{VPN} es crítico, externo y tiene 3 fallos
$\Rightarrow$ \texttt{BLOQUEADO}; \texttt{CORREO} es interno, no crítico y tiene
3 fallos $\Rightarrow$ \texttt{ALERTA}. Como \texttt{BLOQUEADO} $>$
\texttt{ALERTA}, la atención prioritaria es para \texttt{VPN}.

\subsection{Casos de prueba obligatorios}
Cada fila es una prueba independiente de un objeto \texttt{Servidor}. El
resultado obtenido coincide con el esperado en los nueve casos.

\begin{table}[h!]
\centering
\small
\begin{tabular}{@{}llllrrlll@{}}
\toprule
\textbf{Nombre} & \textbf{Origen} & \textbf{Crítico} & \textbf{Total} & \textbf{Fallos} & \textbf{\%} & \textbf{Esperado} & \textbf{Obtenido} \\
\midrule
VPN    & externo & si     & 10 & 3 & 30.0  & BLOQUEADO & BLOQUEADO \\
WEB    & interno & no     & 8  & 4 & 50.0  & ALERTA    & ALERTA \\
CORREO & externo & no     & 8  & 2 & 25.0  & ALERTA    & ALERTA \\
VPN    & interno & si     & 12 & 2 & 16.7  & NORMAL    & NORMAL \\
WEB    & externo & no     & 5  & 5 & 100.0 & BLOQUEADO & BLOQUEADO \\
WEB    & interno & no     & 4  & 5 & N/D   & DATOS INVALIDOS & DATOS INVALIDOS \\
WEB    & interno & no     & 0  & 0 & N/D   & DATOS INVALIDOS & DATOS INVALIDOS \\
CORREO & local   & no     & 10 & 2 & N/D   & DATOS INVALIDOS & DATOS INVALIDOS \\
WEB    & externo & quizás & 10 & 2 & N/D   & DATOS INVALIDOS & DATOS INVALIDOS \\
\bottomrule
\end{tabular}
\caption{Casos de prueba obligatorios.}
\end{table}

\subsection{Casos límite}
\begin{itemize}[nosep]
  \item Entradas con espacios y mayúsculas, como \texttt{"\ EXTERNO\ "} y \texttt{"\ SI\ "}, se interpretan como \texttt{externo} y \texttt{si} gracias a \texttt{trim()} y \texttt{equalsIgnoreCase()}.
  \item Un nombre formado solo por espacios se rechaza como inválido.
  \item Con \texttt{intentosTotales = 0} el porcentaje no se calcula, por lo que no hay división por cero.
  \item Con 5 fallos sobre 5 intentos (extremo incluido) los datos son válidos y el porcentaje es 100.0\%.
\end{itemize}

% =====================================================
\section{Cumplimiento de las restricciones}
\begin{itemize}[nosep]
  \item Un único archivo \texttt{Main.java} con \texttt{class Servidor} y \texttt{public class Main}.
  \item Sin arrays, \texttt{ArrayList}, colecciones ni herencia.
  \item Sin constructores personalizados: los atributos se asignan después de \texttt{new Servidor()}.
  \item Sin modificadores de encapsulación (\texttt{private}, \texttt{public} en atributos, \textit{getters}/\textit{setters}).
  \item Los datos numéricos se introducen como enteros.
\end{itemize}

% =====================================================
\section{Preguntas conceptuales}

\paragraph{1. ¿Por qué dos objetos de la misma clase pueden almacenar valores diferentes en sus atributos?}
Porque la clase es solo la plantilla que define qué atributos existen. Cada
objeto creado con \texttt{new} reserva su propia zona de memoria, con su propia
copia de los atributos, por lo que modificar uno no afecta al otro.

\paragraph{2. ¿Qué diferencia existe entre \texttt{Servidor b = a;} y \texttt{Servidor b = new Servidor();}?}
\texttt{Servidor b = a;} no crea ningún objeto nuevo: copia la referencia, de
modo que \texttt{a} y \texttt{b} apuntan al mismo objeto y un cambio hecho
desde \texttt{b} se ve desde \texttt{a}. \texttt{Servidor b = new Servidor();}
crea un objeto nuevo e independiente, con sus propios atributos.

\paragraph{3. ¿Por qué es importante comprobar \texttt{datosValidos()} antes de calcular el porcentaje?}
Porque con datos inválidos el cálculo no tiene sentido o falla: si
\texttt{intentosTotales} es 0 se produciría una división por cero, y si los
fallidos son negativos o mayores que el total el porcentaje sería incoherente
(negativo o superior al 100\%).

% =====================================================
\section{Conclusiones}
El programa cumple los requisitos del enunciado: separa responsabilidades en
una clase \texttt{Servidor} con métodos propios, valida los datos antes de
clasificarlos, respeta el orden de prioridad de las reglas, trata de forma
segura los datos inválidos y compara los dos servidores mediante condiciones y
no con respuestas fijas. El uso de \texttt{nextLine()} con conversión a entero y
de \texttt{trim()}/\texttt{equalsIgnoreCase()} hace la entrada robusta frente a
espacios y mayúsculas.

\end{document}
